% 第六章：内核测试与可视化
% 负责人：龙宇海

\section{内核测试与可视化} \label{sec:test}

\subsection{测试环境}

\subsubsection{硬件环境}
\begin{itemize}
    \item 处理器：Intel Xeon Gold 6248R (24核心，48线程)
    \item 内存：64GB DDR4-2933
    \item 缓存：L1 32KB(数据) + 32KB(指令)，L2 1MB，L3 35.75MB
    \item 操作系统：Ubuntu 20.04 LTS
\end{itemize}

\subsubsection{软件环境}
\begin{itemize}
    \item 编译器：GCC 9.4.0
    \item 性能分析工具：Intel VTune Profiler
    \item 基准测试：自研DAG任务生成器
\end{itemize}

\subsection{测试方法}

\subsubsection{基准算法}
\begin{itemize}
    \item \textbf{Global-EDF}：全局最早截止时间优先
    \item \textbf{Partitioned-EDF}：分区最早截止时间优先
    \item \textbf{Work-Stealing}：工作窃取算法
    \item \textbf{Cache-Aware}：现有缓存感知调度算法
\end{itemize}

\subsubsection{性能指标}
\begin{itemize}
    \item \textbf{平均响应时间}：任务从到达到完成的时间
    \item \textbf{截止时间违反率}：超过截止时间的任务比例
    \item \textbf{缓存命中率}：L1、L2、L3缓存的命中率
    \item \textbf{系统吞吐量}：单位时间完成的任务数
\end{itemize}

\subsection{测试结果}

\subsubsection{性能对比}
\begin{table}[H]
\centering
\caption{各调度算法性能对比}
\begin{tabular}{|l|c|c|c|c|}
\hline
\textbf{算法} & \textbf{平均响应时间(ms)} & \textbf{截止时间违反率(\%)} & \textbf{L1缓存命中率(\%)} & \textbf{系统吞吐量} \\
\hline
Global-EDF & 45.2 & 12.3 & 78.5 & 1.0 \\
\hline
Partitioned-EDF & 42.1 & 10.8 & 82.1 & 1.08 \\
\hline
Work-Stealing & 38.7 & 9.2 & 84.3 & 1.15 \\
\hline
Cache-Aware & 35.4 & 7.6 & 87.2 & 1.25 \\
\hline
\textbf{Yat-CASched} & \textbf{28.9} & \textbf{4.1} & \textbf{91.7} & \textbf{1.42} \\
\hline
\end{tabular}
\end{table}

\subsubsection{可扩展性测试}
随着核心数量的增加，Yat-CASched相对于基准算法的性能优势更加明显，在24核心系统中性能提升达到42\%。

\subsection{测试分析}

测试结果表明，Yat-CASched在所有关键性能指标上都优于现有算法：
\begin{itemize}
    \item 平均响应时间降低36\%
    \item 截止时间违反率降低66\%
    \item L1缓存命中率提升17\%
    \item 系统吞吐量提升42\%
\end{itemize}
